\question[4]
Das Programm zum Zahlenraten aus dem Unterricht enthält die folgende Funktion.
\begin{lstlisting}
def pruefen():
    tipp = int(eingabe.get())
    if tipp < zahl:
        ausgabe['text'] = 'zu klein'
    elif tipp > zahl:
        ausgabe['text'] = 'zu groß'
    else:
        ausgabe['text'] = 'richtig'
\end{lstlisting}
Das Programm soll zusätzlich in einem Label \texttt{versuche} die Anzahl der Versuche anzeigen. Das Label wurde mit \texttt{text='0'} erzeugt.
Welche Zeilen ergänzen wir in der Funktion?
\begin{solutionbox}{2.5cm}
\begin{lstlisting}
anzahl = int(versuche['text'])
versuche['text'] = anzahl + 1
\end{lstlisting}
\end{solutionbox}
